

\newcommand{\FMMLx}{FMML\textsuperscript{X}}

\newcommand{\chapterSuffixA}{\addtocounter{chapter}{-1}
\renewcommand*{\thechapter}{\arabic{chapter}a}}

\newcommand{\chapterBackToNormal}{\renewcommand*{\thechapter}{\arabic{chapter}}}

\newcommand{\chapterA}[1]{
\addtocounter{chapter}{-1}
\renewcommand*{\thechapter}{\arabic{chapter}a}
\chapter{#1}
\renewcommand*{\thechapter}{\arabic{chapter}}}

\newfloat{lstfloat}{htbp}{lop}
\floatname{lstfloat}{Listing}

\newcommand{\lstRef}[1]{Listing~\ref{#1}}
\newcommand{\conRef}[1]{Con.~\ref{#1}}

\newcounter{constCounter}[section] 
\makeatletter
\newcommand{\constLabel}[1]{%
		\refstepcounter{constCounter}
		\def\@currentlabel{(C \arabic{constCounter})}% Update label
		\raisebox{\f@size pt}\phantomsection
		\label{#1}\hspace{-2.6mm}
		(C \arabic{constCounter})}
\makeatother

\newcounter{defiCounter}[section] 
\makeatletter
\newcommand{\defiLabel}[1]{%
		\refstepcounter{defiCounter}
		\def\@currentlabel{(D \arabic{defiCounter})}% Update label
		\raisebox{\f@size pt}\phantomsection
		\label{#1}\hspace{-2.6mm}
		(D \arabic{defiCounter})}
\makeatother

\newcounter{myCounter}[section] 
\newcounter{mySubCounter}[myCounter] 

\makeatletter
\newcommand{\reqLabel}[1]{%
\myLabel{#1}{Req}}
\makeatother

%%%%%%%%%%%%%%%%%%%%%%%%%%%%%%%%%% Variante 1 %%%%%%%%%%%%%%%%%%%%%%%%%%%%%%%%%%%%%%%
\newcommand{\layerOne}[1]{\chapter{#1}}
\newcommand{\layerOneStar}[1]{\chapter*{#1}}
\newcommand{\layerTwo}[1]{\section{#1}}
\newcommand{\layerThree}[1]{\subsection{#1}}
\newcommand{\layerFour}[1]{\subsubsection{#1}}

\newcommand{\note}[2]{\noindent\fcolorbox{black}{yellow!20}{\begin{minipage}{\dimexpr\linewidth-2\fboxsep-2\fboxrule}\cite{#2}

#1\end{minipage}}}

\makeatletter
	\newcommand{\myLabel}[2]{%
		\refstepcounter{myCounter}
		\def\@currentlabel{#2-\thesubsection.\arabic{myCounter}}% Update label
		\raisebox{\f@size pt}\phantomsection
		\label{#1}\hspace{-2.6mm}
		#2-\thesubsection.\arabic{myCounter}}
\makeatother

\makeatletter
	\newcommand{\myLabelPrefix}[3]{%
		\refstepcounter{#3}
		\def\@currentlabel{#2\arabic{#3}}% Update label
		\raisebox{\f@size pt}\phantomsection
		\label{#1}\hspace{-2.6mm}
		#2\arabic{#3}}
\makeatother

\makeatletter
	\newcommand{\myLabelPrefixNoCount}[3]{%
%		\refstepcounter{#3}
		\def\@currentlabel{#2#3}% Update label
		\raisebox{\f@size pt}\phantomsection
		\label{#1}\hspace{-2.6mm}
		#2#3}
\makeatother

%\makeatletter
%	\newcommand{\myLabelFour}[2]{%
%		\refstepcounter{myCounter}
%		\def\@currentlabel{#2-\thesubsection.\arabic{myCounter}}% Update label
%		\raisebox{\f@size pt}\phantomsection
%		\label{req:#1}
%		#2-\thesubsection.\arabic{myCounter}}
%\makeatother
%
%\newcommand{\kmh}{$kmh^{-1}$}

\makeatletter
	\newcommand{\myLabelB}[2]{%
		\refstepcounter{mySubCounter}
		\def\@currentlabel{#2-\thesubsection.\arabic{myCounter}\alph{mySubCounter}}% Update label
		\raisebox{\f@size pt}\phantomsection
		\label{#1}\hspace{-2.6mm}
		#2-\thesubsection.\arabic{myCounter}\alph{mySubCounter}}
\makeatother

%%%%%%%%%%%%%%%%%%%%%%%%%%%%%%%%%% Variante 2 %%%%%%%%%%%%%%%%%%%%%%%%%%%%%%%%%%%%%%%

%\newcommand{\layerOne}[1]{\part{#1}}
%\newcommand{\layerOneStar}[1]{\part*{#1}}
%\newcommand{\layerTwo}[1]{\chapter{#1}}
%\newcommand{\layerThree}[1]{\section{#1}}

%\makeatletter
%	\newcommand{\myLabel}[2]{%
%		\refstepcounter{myCounter}
%		\def\@currentlabel{#2-\thesubsection.\arabic{myCounter}}% Update label
%		\raisebox{\f@size pt}\phantomsection
%		\label{req:#1}
%		#2-\thesubsection.\arabic{myCounter}}
%\makeatother

%\makeatletter
%	\newcommand{\myLabelB}[2]{%
%		\refstepcounter{mySubCounter}
%		\def\@currentlabel{#2-\thesubsection.\arabic{myCounter}\alph{mySubCounter}}% Update label
%		\raisebox{\f@size pt}\phantomsection
%		\label{req:#1}
%		#2-\thesubsection.\arabic{myCounter}\alph{mySubCounter}}
%\makeatother

%%%%%%%%%%%%%%%%%%%%%%%%%%%%%%%%%% Variante Ende %%%%%%%%%%%%%%%%%%%%%%%%%%%%%%%%%%%%%%%

\makeatletter
  \newcommand{\myRef}[1]{
  \ref{req:#1}}
\makeatother

\makeatletter
  \newcommand{\comment}{
  \hspace{1em} \textit{Comment}}
\makeatother

%%%%%%%%%%%%%%%%%%%%%%%%%%%%%%%%%%%%%%% XOCL %%%%%%%%%%%%%%%%%%%%%%%%%%%%%%%%%%%%%%%%%%%

\definecolor{mygreen}{rgb}{0,0.6,0}
\definecolor{mygray}{rgb}{0.5,0.5,0.5}
\definecolor{mymauve}{rgb}{0.58,0,0.82}
\definecolor{javapurple}{rgb}{0.5,0,0.35} % keywords
\definecolor{main-color}{rgb}{0.6627, 0.7176, 0.7764}
\definecolor{back-color}{rgb}{0.1686, 0.1686, 0.1686}
\definecolor{black-color}{rgb}{0,0,0}
\definecolor{string-color}{rgb}{0.3333, 0.5254, 0.345}
\definecolor{comment-color}{rgb}{0.0, 0.2, 0.0}
\definecolor{key-color}{rgb}{0.8, 0.47, 0.196}
\definecolor{eclipseBlue}{RGB}{42,0.0,255}
\definecolor{ref-class}{RGB}{200,100,100}
\definecolor{eclipsePurple}{RGB}{127,0,85}
\usepackage[T1]{fontenc}

%\renewcommand{\lstlistingname}{Constraint}

\def\myfontsize{\scriptsize}
\lstdefinelanguage{XOCL}{
    basicstyle = {\myfontsize\ttfamily},
    stringstyle = {\color{string-color}},
    comment=[l][\color{comment-color}]{//},%
    keywordstyle = [1]{\myfontsize\ttfamily\bf\color{black-color}},
    keywordstyle = [2]{\myfontsize\ttfamily\color{ref-class}},
    keywordstyle = [3]{\myfontsize\ttfamily\color{teal}},
    keywordstyle = [4]{\myfontsize\ttfamily\color{orange}},    %do,end,grab,or,orelse,case,extends,for,null,data,if,else,in,when,then,let,not,union,letrec,act,self,become,export,probably,and},
    morekeywords = [1]{import,metapackage,and,andthen,or,orelse,not,forAll,context,intern,reason,if,then,else,implies,elseif,end,do,let,in,when,extends,metaclass},
    morekeywords = [2]{Slot,Package,Obj,Boolean,Root,String,Integer,Operation,For,Find,Class,Attribute,Constructor,AbstractOp,Constraint},
    morekeywords = [3]{theClassClass},
    morekeywords = [4]{self,super},  
    literate = {@}{{\textcolor{eclipseBlue}{\ttfamily @}}}1
               {->}{{$\rightarrow$}}2
               {implies}{{$\Rightarrow$}}2
               {subseteq}{{$\subseteq$}}1
               {RefOperation}{{\color{ref-class}Operation}}9
} 

\lstdefinelanguage{XOCLInline}{
    basicstyle = {\myfontsize\ttfamily},
    stringstyle = {\color{string-color}},
    comment=[l][\color{comment-color}]{//},%
    keywordstyle = [1]{\myfontsize\ttfamily\bf\color{black-color}},
    keywordstyle = [2]{\myfontsize\ttfamily\color{ref-class}},
    keywordstyle = [3]{\myfontsize\ttfamily\color{teal}},
    keywordstyle = [4]{\myfontsize\ttfamily\color{orange}},    %do,end,grab,or,orelse,case,extends,for,null,data,if,else,in,when,then,let,not,union,letrec,act,self,become,export,probably,and},
    morekeywords = [1]{and,andthen,or,orelse,not,context,intern,if,then,forAll,else,elseif,end,do,let,in,when,extends,metaclass},
    literate = {@}{{\textcolor{eclipseBlue}{\ttfamily @}}}1
               {->}{{$\rightarrow$}}2
               {implies}{{$\Rightarrow$}}2
               {subseteq}{{$\subseteq$}}1
               {RefOperation}{{\color{ref-class}Operation}}9
} 

\lstset{
  language={XOCL},
  mathescape,
  columns=fixed, % make all characters equal width
  keepspaces=true, % does not ignore spaces to fit width, convert tabs to spaces
  showstringspaces=false, % lets spaces in strings appear as real spaces
  breaklines=true, % wrap lines if they don't fit
  frame=trbl, % draw a frame at the top, right, left and bottom of the listing
  escapeinside={(*}{*)},
  numberstyle=\tiny\ttfamily,
  stepnumber=2,                    % the step between two line-numbers. If it's 1, each line will be numbered
  captionpos=b,                    % sets the caption-position to bottom
  tabsize=2,	                   % sets default tabsize to 2 spaces
	title=\lstname
}

\lstnewenvironment{XOCL}[2]
    {\lstset{language={XOCL}, caption=#1, label=#2}}
    {}

\lstnewenvironment{inlineXOCL}
    {\lstset{language={XOCL},frame=none}}
    {}

 \def\code#1{{\normalfont\lstinline[language={XOCLInline},mathescape,basicstyle=\small\ttfamily]{#1}}}

%%%%%%%%%%%%%%%%%%%%%%%%%%%%%%%%%%

\newcommand{\consoleTable}[3]{
\begin{tabular}{|p{0.8\textwidth}|} \hline
\textbf{In:} #1 \\ \hline
\textbf{Out:} #2 \\ \hline
#3\\ \hline
\end{tabular}
}
